\chapter{Existence of gcds, and B\'ezout's identity}
\label{ch:bezout}

This chapter discharges the one step the informal source leaves open. The source
records the statement

\begin{quote}
  \emph{if $a$, $b$ are naturals, there exist $x, y$ integers such that
  $x a + y b = \gcd(a,b)$}
\end{quote}

and then says only \texttt{proof: TODO}. The argument below is therefore the
project's own; there is no source proof to quote. We take the Euclidean-algorithm
route --- strong induction along $b \bmod a$ --- rather than the
least-positive-element-of $\{xa + yb\}$ route, because it needs no well-ordering
setup and produces the gcd and the coefficients in a single recursion.

\section{Coercion discipline}

The source writes $x a + y b = \gcd(a,b)$ without comment, but $a$, $b$ and
$\gcd(a,b)$ are naturals while $x$ and $y$ are integers. Every statement in this
project fixes the convention once, here, and every downstream chapter inherits
it:

\begin{itemize}
  \item \textbf{The predicates $\Prm{\cdot}$, $\Irr{\cdot}$, $\Gcd{\cdot}{\cdot}{\cdot}$ and
    every divisibility hypothesis live in $\NN$.} They are never restated over
    $\ZZ$.
  \item \textbf{$\ZZ$ appears only for the two B\'ezout coefficients}, and only
    inside the equation they satisfy. The equation is written with the inclusion
    $\iota : \NN \to \ZZ$ applied to each natural: $x\iota(a) + y\iota(b) = \iota(d)$.
  \item \textbf{Crossing back is by the divisibility bridge}: $\iota(m) \mid \iota(n)$
    in $\ZZ$ implies $m \mid n$ in $\NN$. This is the only bridge
    Chapter~\ref{ch:irreducible} needs, and it is used exactly once.
  \item \textbf{Truncated subtraction is never used inside a cast.} Wherever the
    Euclidean step wants $b - a q$, we instead use the \emph{addition} identity
    $a \cdot (b \operatorname{div} a) + (b \bmod a) = b$ and transport that
    equation into $\ZZ$, where subtraction is total. This avoids the defect the
    source's proof of Lemma~\ref{lem:dvd_lt_of_ne} ran into.
\end{itemize}

\section{Uniqueness}

\begin{lemma}[a gcd is unique]
  \leanok
  \label{lem:is_gcd_unique}
  \lean{EuclidPrimes.IsGcd.unique}
  \uses{def:is_gcd}
  If $\Gcd{a}{b}{d}$ and $\Gcd{a}{b}{d'}$ then $d = d'$.
\end{lemma}
\begin{proof}
  \uses{def:is_gcd}
  \leanok
  $d'$ divides $a$ and divides $b$, so the universal clause of $\Gcd{a}{b}{d}$
  applied to $c := d'$ gives $d' \mid d$. Symmetrically $d \mid d'$. Two natural
  numbers that divide each other are equal.
\end{proof}

This is what lets us speak of ``the'' gcd, and it is what will let us upgrade
B\'ezout from the particular $d$ the construction produces to an arbitrary
witness.

\section{The Euclidean step}

\begin{lemma}[the Euclidean reduction preserves gcds]
  \leanok
  \label{lem:is_gcd_of_mod}
  \lean{EuclidPrimes.isGcd_of_mod}
  \uses{def:is_gcd}
  Let $a, b, d \in \NN$. If $\Gcd{b \bmod a}{a}{d}$ then $\Gcd{a}{b}{d}$.

  No positivity hypothesis is needed. The Euclidean recursion only ever
  applies this lemma with $0 < a$, but the statement holds for every $a$: at
  $a = 0$ we have $b \bmod 0 = b$, and the hypothesis $\Gcd{b}{0}{d}$ gives
  $\Gcd{0}{b}{d}$ by symmetry of the two divisibility clauses.
\end{lemma}
\begin{proof}
  \uses{def:is_gcd}
  \leanok
  Write $q := b \operatorname{div} a$ and $r := b \bmod a$, so that
  \begin{equation}\label{eq:divmod}
    a q + r = b .
  \end{equation}
  (Equation~\eqref{eq:divmod} holds for $a = 0$ as well, with $q = 0$ and
  $r = b$, so nothing below needs $0 < a$.)
  We check the three clauses of $\Gcd{a}{b}{d}$, given $d \mid r$, $d \mid a$,
  and the universal property of $d$ for the pair $(r, a)$.
  \begin{enumerate}
    \item $d \mid a$: this is the second clause of the hypothesis.
    \item $d \mid b$: from $d \mid a$ we get $d \mid a q$, and $d \mid r$, so
      $d$ divides the sum $a q + r = b$ by \eqref{eq:divmod}.
    \item Let $c \in \NN$ with $c \mid a$ and $c \mid b$. A divisor of the
      modulus transports divisibility across the remainder: for any $c \mid a$,
      $c \mid b \bmod a$ holds exactly when $c \mid b$. (Read left to right this
      is clause 2 above with $c := d$; read right to left it gives $c \mid r$
      here. Both directions follow from \eqref{eq:divmod} and $c \mid aq$
      without ever forming a difference in $\NN$ --- see the coercion
      discipline: no truncated subtraction anywhere.) So $c \mid r$, and the
      universal clause of $\Gcd{r}{a}{d}$, applied to this $c$ with $c \mid r$
      and $c \mid a$, gives $c \mid d$.
  \end{enumerate}
\end{proof}

\section{The main construction}

\begin{theorem}[simultaneous existence of a gcd and of B\'ezout coefficients]
  \leanok
  \label{thm:exists_is_gcd_bezout}
  \lean{EuclidPrimes.exists_isGcd_bezout}
  \uses{def:is_gcd, lem:is_gcd_of_mod}
  For all $a, b \in \NN$ there exists $d \in \NN$ such that $\Gcd{a}{b}{d}$ and
  there exist $x, y \in \ZZ$ with $x\iota(a) + y\iota(b) = \iota(d)$.
\end{theorem}
\begin{proof}
  \uses{def:is_gcd, lem:is_gcd_of_mod}
  \leanok
  The two conclusions are proved \emph{together}, in one induction. Separating
  them would force the induction to be run twice, and the second run would have
  no handle on \emph{which} $d$ the first one produced.

  We argue by strong induction on $a$, with $b$ universally quantified inside
  the induction hypothesis. So fix $a$ and assume:
  \[
    \text{for every } a' < a \text{ and every } b' \in \NN,\ \exists d,\ \Gcd{a'}{b'}{d}
    \ \text{and}\ \exists x, y \in \ZZ,\ x\iota(a') + y\iota(b') = \iota(d).
  \]

  \paragraph{Case $a = 0$.} Take $d := b$. Then $\Gcd{0}{b}{b}$ holds: $b \mid 0$
  since everything divides $0$; $b \mid b$; and if $c \mid 0$ and $c \mid b$ then
  $c \mid b$. For the coefficients take $x := 0$, $y := 1$, giving
  $0 \cdot \iota(0) + 1 \cdot \iota(b) = \iota(b)$.

  \paragraph{Case $0 < a$.} Put $q := b \operatorname{div} a$ and
  $r := b \bmod a$. Since $0 < a$ we have $r < a$, so the induction hypothesis
  applies to the pair $(r, a)$ --- note the arguments are \emph{swapped}, which
  is exactly what makes the recursion decrease. It yields $d \in \NN$ with
  $\Gcd{r}{a}{d}$ and $x, y \in \ZZ$ with
  \begin{equation}\label{eq:ih}
    x \iota(r) + y \iota(a) = \iota(d).
  \end{equation}

  By Lemma~\ref{lem:is_gcd_of_mod} we get $\Gcd{a}{b}{d}$, which is the first
  conclusion. (That lemma carries no positivity hypothesis; $0 < a$ is needed
  here only to know $r < a$, which is what makes the recursion decrease.)

  For the second, transport the division identity $a q + r = b$ into $\ZZ$:
  \[
    \iota(a)\,\iota(q) + \iota(r) = \iota(b),
    \qquad\text{hence}\qquad
    \iota(r) = \iota(b) - \iota(a)\,\iota(q).
  \]
  Substituting into \eqref{eq:ih}:
  \[
    \iota(d) = x\bigl(\iota(b) - \iota(a)\iota(q)\bigr) + y\,\iota(a)
             = \bigl(y - x\,\iota(q)\bigr)\iota(a) + x\,\iota(b).
  \]
  So the coefficients for the pair $(a, b)$ are
  \[
    x' := y - x\,\iota(q), \qquad y' := x ,
  \]
  and $x'\iota(a) + y'\iota(b) = \iota(d)$ as required. The manipulation is pure
  ring algebra in $\ZZ$; no subtraction ever occurs in $\NN$.
\end{proof}

\section{The two consequences used downstream}

\begin{corollary}[every pair has a gcd]
  \leanok
  \label{cor:exists_is_gcd}
  \lean{EuclidPrimes.exists_isGcd}
  \uses{def:is_gcd, thm:exists_is_gcd_bezout}
  For all $a, b \in \NN$ there exists $d \in \NN$ with $\Gcd{a}{b}{d}$.
\end{corollary}
\begin{proof}
  \uses{thm:exists_is_gcd_bezout}
  \leanok
  Discard the B\'ezout half of Theorem~\ref{thm:exists_is_gcd_bezout}.
\end{proof}

\begin{lemma}[B\'ezout, in the form the source states it]
  \leanok
  \label{lem:bezout}
  \lean{EuclidPrimes.bezout}
  \uses{def:is_gcd, thm:exists_is_gcd_bezout, lem:is_gcd_unique}
  % SOURCE: [euclid-notes], the first `lemma:` (read from references/euclid-infinitude-of-primes.md)
  % SOURCE QUOTE: "lemma: if a b are naturals, there exists x,y integers such that x*a + y * b = gcd(a,b)
  % proof: TODO"
  \textit{Source: euclid-infinitude-of-primes.md, first \texttt{lemma:} (stated there with \texttt{proof: TODO}).}
  Let $a, b, d \in \NN$ with $\Gcd{a}{b}{d}$. Then there exist $x, y \in \ZZ$
  with
  \[
    x \cdot \iota(a) + y \cdot \iota(b) = \iota(d) .
  \]
\end{lemma}
\begin{proof}
  \uses{thm:exists_is_gcd_bezout, lem:is_gcd_unique}
  \leanok
  Theorem~\ref{thm:exists_is_gcd_bezout} supplies some $d_0$ with
  $\Gcd{a}{b}{d_0}$ together with $x, y \in \ZZ$ satisfying
  $x\iota(a) + y\iota(b) = \iota(d_0)$. By Lemma~\ref{lem:is_gcd_unique} applied
  to $d_0$ and $d$, we get $d_0 = d$, so the same $x, y$ work for $d$.
\end{proof}
